\scp{\tsc{Sumba-Havu}}{03-03}
While several modern anthropological studies have been done across
much of Sumba, linguistic descriptions are scarce.
\citet{Klamer1998a}, on \ili{Kambera}, is the one full modern grammar available.
Two other modern linguistic descriptions have appeared in recent
years: \citet{GhanggoAte2018} describes reduplication in \ili{Kodi},
and \citet{Verdizade2019a} describes the phonology
and aspects of the morphosyntax of \ili{Laboya}.
\citet{Onvlee1984} is a valuable source of data,
though it has transcriptional issues (see \srf{03-03-02}).
Three sources which provide perspectives on the linguistic context
and speech varieties in Sumba are: \citet{Asplund2010},
\citet{GhanggoAte2020}, and \citet{Lovestrand2021}.

\citet{Grimes2010a} shows that while \ili{Havu}
and \ili{Dhao} are relatively close in their lexicons
and historical phonologies, they are significantly different
in their syntax and typology.
Both \ili{Havu} and \ili{Dhao} have six-vowel systems,
with penultimate schwa /ə/ triggering phonetic lengthening of
the following consonant (see \srf{03-03-01}).
Both have four implosives.
Phonetically long vowels are phonemically a sequence of two like vowels.
\ili{Dhao} has an extra grade of obstruents not found in \ili{Havu}.
\ili{Havu} uses a template-inserted /e/ as its antepenultimate vowel,
whereas \ili{Dhao} uses /a/.
In both languages, V-initial verb roots are inflected for person and number.
\ili{Havu} is VSO and ergative marked.
\ili{Dhao} is SVO and basically nominative-accusative, probably due
to intense contact and bilingualism over time with the nearby languages of western Rote.
\citet{WalkerA1982,Grimes2008,GrimesRanohAplugi2008}, and
\citet{Balukh2020} are important sources of data for these languages.

Speakers of the five dialects of \ili{Havu} on the Sabu mainland claim
to understand each other,
but often with difficulty.
We separate \ili{Raijua} as a separate language,
based on different ergative markers (\ili{Raijua} \it{la},
versus \ili{Seba} \it{ri}, and \ili{Dimu} \it{ro/rovi}),
differences in several pronouns,
and other differences in the lexicon and phonology (see \citealp{Grimes2010a}),
all combining for consistent reports of blocked intelligibility.
\citet{Blust2012a} expands on historical-comparative issues in
\tsc{\ili{Havu}-Dhao} listed in \citet{Grimes2010a}.